%%%%%%%%%%%%%%%%%%%%%%%%%%%%%%%%%%%%%%%%%%%%%%%%%%%%%%%%%%%%
% GLOBAL PREAMBLE BY ENOCH YU
%%%%%%%%%%%%%%%%%%%%%%%%%%%%%%%%%%%%%%%%%%%%%%%%%%%%%%%%%%%%

%%%%%%%%%%%%%%% PACKAGES %%%%%%%%%%%%%%%
% Core
\usepackage[T1]{fontenc}
\usepackage[utf8]{inputenc}
\usepackage{palatino}
\usepackage{mathpazo}

% Math
\usepackage{amsmath}
\usepackage{amssymb}
\usepackage{amsthm}
\usepackage{mathtools}

% Page Layout
\usepackage[margin=1in, headsep=7pt, headheight=14pt]{geometry}
\usepackage{lastpage}
\usepackage{fancyhdr}
\fancyhead[L]{Enoch's Notebook: Linear Algebra}
\fancyhead[R]{Enoch Yu}
\pagestyle{fancy}
\usepackage{parskip}
\usepackage{microtype}
\setlength{\headsep}{7pt}
\setlength{\headheight}{14pt}

% Sections
\usepackage{titlesec}
\titleformat{\section}
  {\normalfont\Large\bfseries\color{BrickRed}}
  {\S\! \thesection}
  {0.5em}
  {}

\usepackage[Lenny]{fncychap}
\ChNameVar{\fontsize{14}{16}\usefont{OT1}{phv}{m}{n}\selectfont}
\ChNumVar{\fontsize{60}{62}\usefont{OT1}{ptm}{m}{n}\selectfont}
\ChTitleVar{\Huge\bfseries\rm}
\ChRuleWidth{1pt}
\renewcommand{\chaptername}{Note}


% Tikz
\usepackage{tikz}
\usetikzlibrary{calc, angles, shapes.geometric, positioning, arrows.meta}
\usepackage{tkz-euclide}
\usepackage{pgfplots}
\pgfplotsset{compat=1.9}
\usepackage[font=small,labelfont=bf]{caption}

% Miscellaneous
\usepackage{enumitem}
\usepackage[dvipsnames]{xcolor}
\usepackage[normalem]{ulem}
\usepackage{url}
\usepackage{hyperref}
\hypersetup{
    colorlinks=true,
    linkcolor=black,
    filecolor=magenta,      
    urlcolor=blue,
    %urlcolor=cyan,
    pdfpagemode=FullScreen,
}
\usepackage{graphicx}
\graphicspath{ {./images/} }
\usepackage{tabularx}
\usepackage{arydshln}
\usepackage{float}
\usepackage{pdfpages}
\usepackage{subfiles}


%%%%%%%%%%%%%%% CUSTOM %%%%%%%%%%%%%%%
% Theorems
\newtheorem{theorem}{Theorem}
\newtheorem{lemma}[theorem]{Lemma}
\newtheorem*{lemma*}{Lemma}
\newtheorem{corollary}[theorem]{Corollary}

% Environments
\newenvironment{solution}
    {\begin{trivlist}\item[]{\bf Solution:}}
    {\qed \end{trivlist}}
\newenvironment{conclusion}
    {\begin{trivlist}\item[]{\bf Conclusion: }\color{BrickRed}}
    {\end{trivlist}}
\newenvironment{ThoughtProcess}
    {\begin{trivlist}\item[]{\bf Thought Process: }\color{MidnightBlue}}
    {\end{trivlist}}

% Boxes
\usepackage{tcolorbox}
\tcbuselibrary{theorems}
\newtcbtheorem[number within=section]{theobox}{Theorem}%
{%
    colback=black!5,
    colframe=black,
    fonttitle=\bfseries,
    sharp corners
}{theo}
\newtcbtheorem[number within=section]{probbox}{Exercise}%
{%
    colback=white,
    colframe=red!40!black,
    fonttitle=\bfseries,
    sharp corners
}{prob}
\newtcbtheorem[number within=section]{corbox}{Corollary}%
{%
    colback=white,
    colframe=blue!40!black,
    fonttitle=\bfseries,
    sharp corners
}{cor}

\usepackage[framemethod=TikZ]{mdframed}
% Definition
\newenvironment{keydef}
{\begin{mdframed}[
    linecolor=Sepia, 
    linewidth=1pt,
    skipabove=\medskipamount,
    skipbelow=\medskipamount,
    backgroundcolor=yellow!10,
    innertopmargin=10pt,
    innerbottommargin=10pt,
    innerleftmargin=10pt,
    innerrightmargin=10pt
]}
{\end{mdframed}}

% Example
\NewDocumentEnvironment{example}{o}
{
 \IfValueT{#1}{\vspace{11pt}}
 \begin{mdframed}[
    linecolor=CadetBlue,
    linewidth=1pt,
    skipabove=\medskipamount,
    skipbelow=\medskipamount,
    backgroundcolor=cyan!5,
    innertopmargin=10pt,
    innerbottommargin=10pt,
    innerleftmargin=10pt,
    innerrightmargin=10pt
 ]
 \IfValueT{#1}{%
    \vspace{-11pt}%
    \noindent
    \begin{tikzpicture}[remember picture, overlay]
        \node[
            draw=CadetBlue,
            line width=1pt,
            fill=cyan!5,
            text=Sepia,
            font=\bfseries\small,
            anchor=west,
            xshift=2pt
        ] {\,\strut #1\,};
    \end{tikzpicture}\par
    \vspace{5pt}%
 }%
}
{\end{mdframed}}

% Thought Process
\newenvironment{thoughts}
{\begin{mdframed}[
    topline=false, 
    bottomline=false, 
    rightline=false,
    linecolor=RoyalPurple, 
    linewidth=2pt,
    skipabove=\medskipamount, 
    skipbelow=\medskipamount,
    fontcolor=RoyalPurple,
    font=\small\itshape
]}
{\end{mdframed}}

% Tip/Side note
\newenvironment{tip}
{\begin{mdframed}[
    topline=false, 
    bottomline=false, 
    rightline=false,
    linecolor=BrickRed, 
    linewidth=2pt,
    skipabove=\medskipamount, 
    skipbelow=\medskipamount,
    fontcolor=BrickRed,
    font=\small\itshape
]}
{\end{mdframed}}


% Jupyter cell like block
\usepackage{listings}
\definecolor{nb@bg}{RGB}{231,234,234}
\definecolor{nb@kw}{RGB}{0,64,153}
\definecolor{nb@str}{RGB}{163,21,21}
\definecolor{nb@cmt}{RGB}{100,100,100}

\lstdefinestyle{jin}{
    language=Python,
    backgroundcolor=\color{nb@bg},
    basicstyle=\ttfamily\small,
    keywordstyle=\color{nb@kw}\bfseries,
    stringstyle=\color{nb@str},
    commentstyle=\color{nb@cmt}\itshape,
    numberstyle=\tiny\color{gray},
    numbers=left, numbersep=8pt,
    breaklines=true,
    frame=single,
    rulecolor=\color{gray!50},
    showstringspaces=false,
    tabsize=4, captionpos=b,
    xleftmargin=1.2em,
    framexleftmargin=1.2em,
    belowskip=0pt,
    morekeywords={
        Sequential,Dense,model,compile,fit,predict,
        Adam,BinaryCrossentropy,SparseCategoricalCrossentropy,
        MeanSquaredError,from_logits,activation,units,
        optimizer,loss,epochs,learning_rate,
        tf,keras,tensorflow,StandardScaler,
        XGBClassifier,XGBRegressor,np,plt
    }
}
\lstdefinestyle{jout}{
    basicstyle=\ttfamily\small,
    breaklines=true,
    showstringspaces=false,
    tabsize=4, captionpos=b,
    aboveskip=4pt,
    xleftmargin=1.2em,
}

\newcounter{jupytercell}
\newenvironment{jupyter}{
    \stepcounter{jupytercell}
    \begin{mdframed}[
        backgroundcolor=white,
        linecolor=Sepia!80,
        linewidth=1pt,
        innerleftmargin=10pt,
        innerrightmargin=10pt,
        innertopmargin=28pt,
        innerbottommargin=0pt,
        skipabove=15pt,
        skipbelow=15pt
    ]
    \vspace{-28pt}%
    \noindent\hspace{-11pt}%
    \begin{tikzpicture}[remember picture, overlay]
        \node[
            draw=Sepia!80,
            line width=1pt,
            fill=white,
            text=Sepia,
            font=\bfseries\small,
            anchor=west,
            xshift=10pt,
        ] {\,\strut Cell \thejupytercell\,};
    \end{tikzpicture}\par
    \vspace{8pt}%
}{
    \end{mdframed}
}





\title{Enoch's Notebook: Linear Algebra}
\author{Enoch Yu}
\date{\today}


%%%%%%%%%%%%%%% END %%%%%%%%%%%%%%%



